\chapter{Systèmes d'équations algébriques linéaires}

\section{Introduction}

Le traitement des problèmes scientifiques conduit très souvent à des systèmes
d'équations algébriques linéaires à plusieurs inconnues. Lorsque le nombre
d'équations n'est pas élevé (par exemple moins de 3), on peut trouver les
solutions du système d'équations par des méthodes analytiques connues
(éliminations, substitution). Mais lorsque ce nombre devient élevé ou lorsque
les coefficients des équations sont variables, il est utile de développer des
codes de calcul numériques qui aideront l'ordinateur à calculer les solutions.

L'objet de ce chapitre est de présenter les schémas numériques pour la
résolution des équations algébriques linéaires : les méthodes directes
(élimination de Gauss et de Jordan, racine carrée de Cholesky, décomposition
triangulaire, fragmentation) et les méthodes itératives (de Jacobi et de
Gauss-Seidel).

\section{Méthodes directes pour les systèmes d'équations linéaires}

\subsection{Méthodes d'élimination de Gauss et de Jordan}

\subsubsection{Principe de la méthode}

Soit le système algébrique de la forme :
\[
  AX = b
\]
où $A = (a_{ij})$, $X = (x_i)$ et $b = (b_i)$.

Ce système d'équations est de la forme
\[
  \left\{
  \begin{aligned}
    a_{11}x_1 + a_{12}x_2 + \dots + a_{1n}x_n &= b_1\\
    a_{21}x_1 + a_{22}x_2 + \dots + a_{2n}x_n &= b_2\\
    &\;\;\vdots\\
    a_{n1}x_1 + a_{n2}x_2 + \dots + a_{nn}x_n &= b_n
  \end{aligned}
  \right.
\]

La méthode d'élimination de Gauss est basée sur le principe suivant. Les
inconnues sont éliminées successivement en exprimant à partir d'une des
équations une inconnue en fonction des autres puis en substituant dans le
reste des équations (voir classes du secondaire). On part ainsi d'un système de
$n$ équations à $n$ inconnues à un système de $n-1$ équations à $n-1$
inconnues. Le processus se poursuit jusqu'à l'obtention d'une équation à une
inconnue. On trouve alors cette dernière inconnue et par un processus inverse,
les autres inconnues sont calculées. À chaque étape, l'équation qui permet
d'exprimer une inconnue en fonction des autres est appelée \emph{équation
pivotale}. Pour éviter les erreurs d'arrondis, l'équation pivotale doit être
celle qui possède le plus grand coefficient de l'inconnue à éliminer.

La procédure de cette méthode peut alors se formuler de la manière suivante :

Pour une ligne n\textsuperscript{o}~$i$, notée $L_i$, on fait l'opération
\[
  L_i \leftarrow L_i - \frac{a_{i1}}{a_{11}}L_1, \qquad i = 2 \text{ jusqu'à } n.
\]

Pour $i$ allant de $2$ à $n$, on obtient alors un nouveau système de la forme
\[
  \left\{
  \begin{aligned}
    a_{11}x_1 + a_{12}x_2 + \dots + a_{1n}x_n &= b_1\\
    0 + a'_{22}x_2 + \dots + a'_{2n}x_n &= b'_2\\
    0 + a'_{32}x_2 + \dots + a'_{3n}x_n &= b'_3\\
    &\;\;\vdots\\
    0 + a'_{n2}x_2 + \dots + a'_{nn}x_n &= b'_n
  \end{aligned}
  \right.
\]
Les coefficients $a'_{ij}$ et $b'_i$ sont donnés par
\[
  a'_{ij} = a_{ij} - \frac{a_{i1}}{a_{11}}a_{1j}, \qquad
  b'_i = b_i - \frac{a_{i1}}{a_{11}}b_1 .
\]
On continue le processus en considérant le système constitué des lignes $2$ à
$n$ et ainsi de suite jusqu'au système constitué des lignes $n-1$ et $n$.

\subsubsection{Algorithme}

Ainsi pour une étape $k$, c'est-à-dire en considérant le système constitué des
lignes $k$ à $n$, l'algorithme suivant est utilisé :

\begin{algo}
  \For{$l$ allant de $k+1$ à $n$}
    \State $b_l \leftarrow b_l - \dfrac{a_{lk}}{a_{kk}}\,b_k$
    \For{$j$ allant de $k+1$ à $n$}
      \State $a_{lj} \leftarrow a_{lj} - \dfrac{a_{lk}}{a_{kk}}\,a_{kj}$
    \EndFor
  \EndFor
\end{algo}

On peut ainsi en déduire l'algorithme de triangularisation du système
d'équations en faisant varier $k$ de $1$ à $n-1$. Le système triangulaire se
présente ainsi
\[
  \left\{
  \begin{array}{l@{\;}l@{\;}l@{\;}c@{\;}l}
    a_{11}x_1 + & a_{12}x_2 + & \dots\dots\dots\dots\dots\dots & + a_{1n}x_n &= b_1\\
    0 + & a_{22}x_2 + & \dots\dots\dots\dots\dots\dots & + a_{2n}x_n &= b_2\\
    0 + & 0 & + a_{33}x_3 \dots\dots\dots\dots & + a_{3n}x_n &= b_3\\
    \vdots & & & & \\
    0 + & 0 + & \dots\dots\dots\dots\dots\dots + 0 & + a_{nn}x_n &= b_n
  \end{array}
  \right.
\]

\NB Il faut noter que les coefficients $a_{ij}$ et $b_k$ du système
triangulaire ne sont pas ceux du départ.

Le système algébrique triangulaire peut se résoudre de la manière suivante :

La dernière équation donne
\[
  x_n = \frac{b_n}{a_{nn}} .
\]
Pour l'inconnue $x_k$, on a
\[
  x_k = \bigl(b_k - a_{k,k+1}x_{k+1} - a_{k,k+2}x_{k+2} - \dots - a_{k,n}x_n\bigr)\big/ a_{kk}
\]
ou sous forme condensée
\[
  x_k = \Bigl(b_k - \sum_{l=k+1}^{n} a_{k,l}\,x_l\Bigr)\Big/ a_{kk} .
\]

\begin{exercice}[Exercices]
\begin{enumerate}
  \item Écrire l'algorithme de résolution du système algébrique triangulaire.
  \item Écrire l'algorithme qui permet d'interchanger les lignes des systèmes
        algébriques dû à la règle de pivot maximal.
  \item Écrire l'algorithme complet de la méthode de Gauss (triangularisation
        et résolution du système triangulaire) en prenant soin de faire jouer
        la règle du pivot maximal.
\end{enumerate}
\end{exercice}

\subsubsection{Méthode de Jordan}

Elle est analogue à la méthode de Gauss. Mais ici, on met aussi à zéro les
termes qui se trouvent au-dessus de la diagonale principale. On obtient alors
un système diagonal.

\subsection{Méthodes par décomposition de la matrice}

\subsubsection{Méthode de la décomposition triangulaire}

\paragraph{a- Principe}

Soit le système
\[
  AX = b .
\]
On peut de manière unique décomposer la matrice $A$ en deux matrices
triangulaires dont l'une est triangulaire inférieure avec les éléments
quelconques sur la diagonale et l'autre triangulaire supérieure avec tous les
éléments égaux à 1 sur la diagonale.

Soit $L$ et $S$ ces deux matrices triangulaires, on a
\[
  A = LS .
\]
Le système d'équations devient
\[
  SX = L^{-1}b .
\]
Ce dernier système est un système triangulaire que l'on peut résoudre par le
schéma précédent.

\paragraph{Décomposition de la matrice}

On peut écrire $L$ et $S$ sous la forme
\[
  L = \begin{pmatrix}
    l_{11} & 0 & 0 & \cdots & 0\\
    l_{21} & l_{22} & 0 & \cdots & 0\\
    \vdots & & \ddots & & \vdots\\
    \vdots & & & \ddots & 0\\
    l_{n1} & l_{n2} & l_{n3} & \cdots & l_{nn}
  \end{pmatrix},
  \qquad
  S = \begin{pmatrix}
    1 & s_{12} & s_{13} & \cdots & s_{1n}\\
    0 & 1 & s_{23} & \cdots & s_{2n}\\
    \vdots & & 1 & & s_{3n}\\
    \vdots & & & \ddots & \vdots\\
    0 & 0 & 0 & \cdots & 1
  \end{pmatrix}.
\]
L'équation $LS = A$ entraîne par identification les relations suivantes
\[
\begin{array}{ll}
  \left\{\begin{aligned}
    a_{j1} &= l_{j1}\\
    a_{1j} &= l_{11}s_{1j}
  \end{aligned}\right.
  &
  \begin{aligned}
    &1 \le j \le n\\
    &2 \le j \le n
  \end{aligned}
  \\[1.2em]
  \left\{\begin{aligned}
    a_{j2} &= l_{j1}s_{12} + l_{j2}\\
    a_{2j} &= l_{21}s_{1j} + l_{22}s_{2j}
  \end{aligned}\right.
  &
  \begin{aligned}
    &2 \le j \le n\\
    &3 \le j \le n
  \end{aligned}
  \\[1.2em]
  \left\{\begin{aligned}
    a_{j3} &= l_{j1}s_{13} + l_{j2}s_{23} + l_{j3}\\
    a_{3j} &= l_{31}s_{1j} + l_{32}s_{2j} + l_{33}s_{3j}
  \end{aligned}\right.
  &
  \begin{aligned}
    &3 \le j \le n\\
    &4 \le j \le n
  \end{aligned}
  \\
  \quad\vdots & \\
  \left\{\begin{aligned}
    a_{j,n-1} &= l_{j1}s_{1,n-1} + l_{j2}s_{2,n-1} + \dots + l_{j,n-2}s_{n-2,n-1} + l_{j,n-1}\\
    a_{n-1,j} &= l_{n-1,1}s_{1j} + l_{n-1,2}s_{2j} + \dots + l_{n-1,n-1}s_{n-1,j}
  \end{aligned}\right.
  &
  \begin{aligned}
    &n-1 \le j \le n\\
    &j = n
  \end{aligned}
  \\[1.2em]
  \;a_{nn} = l_{n1}s_{1n} + l_{n2}s_{2n} + \dots + l_{nn} &
\end{array}
\]

Avec ces équations on obtient les éléments de la première colonne de $L$ puis
ceux de la première ligne de $S$, ensuite ceux de la 2\textsuperscript{e}
colonne de $L$ et ceux de la 2\textsuperscript{e} ligne de $S$ et ainsi de
suite. On obtient alors les éléments suivants qui permettent de calculer les
éléments $l_{ij}$ et $s_{ij}$. On a :
\[
\begin{aligned}
  &\left\{\begin{aligned}
    l_{j1} &= a_{j1}\\
    s_{1j} &= \frac{a_{1j}}{l_{11}}
  \end{aligned}\right.
  \qquad
  \left\{\begin{aligned}
    l_{j2} &= a_{j2} - l_{j1}s_{12}\\
    s_{2j} &= \frac{1}{l_{22}}\bigl(a_{2j} - l_{21}s_{1j}\bigr)
  \end{aligned}\right.
  \\[0.6em]
  &\left\{\begin{aligned}
    l_{j3} &= a_{j3} - l_{j1}s_{13} - l_{j2}s_{23}\\
    s_{3j} &= \frac{1}{l_{33}}\bigl(a_{3j} - l_{31}s_{1j} - l_{32}s_{2j}\bigr)
  \end{aligned}\right.
  \\
  &\quad\vdots\\
  &\left\{\begin{aligned}
    l_{j,n-1} &= a_{j,n-1} - \sum_{k=1}^{n-2} l_{jk}\,s_{k,n-1}\\
    s_{n-1,j} &= \frac{1}{l_{n-1,n-1}}\Bigl[a_{n-1,j} - \sum_{k=1}^{n-2} l_{n-1,k}\,s_{kj}\Bigr]
  \end{aligned}\right.
  \\[0.6em]
  &\; l_{nn} = a_{nn} - \sum_{k=1}^{n-1} l_{nk}\,s_{kn}
\end{aligned}
\]

\begin{exercice}
\begin{enumerate}
  \item Écrire l'algorithme permettant de calculer les $l_{ij}$ et $s_{ij}$.
  \item Trouver dans la littérature l'algorithme d'inversion d'une matrice et
        de calcul des éléments du produit $L^{-1}b$.
\end{enumerate}
\end{exercice}

\subsubsection{Méthode de la racine carrée de Cholesky}

Une matrice $A$ symétrique peut être mise sous la forme d'un produit de deux
matrices $M$ et $M'$ dont l'une est la transposée de l'autre. On a ainsi une
seule matrice triangulaire supérieure $M$ à calculer. Ceci simplifie les
calculs et diminue l'encombrement mémoire de la machine. En utilisant le
procédé précédent on établit que les éléments de $M$ sont définis par :
\[
\begin{aligned}
  m_{11} &= \sqrt{a_{11}}\\
  m_{1j} &= \frac{a_{1j}}{m_{11}}\\
  m_{22} &= \sqrt{a_{22} - \sum_{k=1}^{1} m_{k2}^2} = \sqrt{a_{22} - m_{12}^2}\\
  m_{2j} &= \frac{a_{2j} - \sum_{k=1}^{1} m_{k2}m_{kj}}{m_{22}}
          = \frac{a_{2j} - m_{12}m_{1j}}{m_{22}}\\
  m_{pp} &= \sqrt{a_{pp} - \sum_{k=1}^{p-1} m_{kp}^2}
          && \text{pour } p \text{ allant de } 2 \text{ à } n\\
  m_{pj} &= \frac{a_{pj} - \sum_{k=1}^{p-1} m_{kp}m_{kj}}{m_{pp}}
          && \text{pour } p \text{ allant de } 2 \text{ à } n
             \text{ et } j \text{ allant de } p+1 \text{ à } n\\
  m_{nn} &= \sqrt{a_{nn} - m_{1n}^2 - m_{2n}^2 - \dots - m_{n-1,n}^2}
\end{aligned}
\]

\begin{exercice}
Écrire l'algorithme de calcul des coefficients $m_{pp}$ et $m_{pj}$.
\end{exercice}

\subsubsection{Fragmentation des systèmes volumineux}

Quand $n$ est très grand, il n'y a pas assez de mémoire pour stocker les
$n(n+1)$ coefficients du système et mettre en place le programme de calcul. On
fragmente alors le système d'équations en décomposant la matrice du système en
sous-matrices.

Par exemple, on mettra le système
\[
  AX = b
\]
sous la forme
\[
  {\renewcommand{\arraystretch}{1.5}
  \left(\begin{array}{c:c}
    A_1 & A_2\\ \hdashline
    A_3 & A_4
  \end{array}\right)
  \left(\begin{array}{c}
    X_1\\ \hdashline
    X_2
  \end{array}\right)
  =
  \left(\begin{array}{c}
    B_1\\ \hdashline
    B_2
  \end{array}\right)}.
\]
Les matrices $A_1$ et $A_2$ ne sont pas nécessairement des matrices carrées et
la décomposition par des droites en pointillés est telle que les
multiplications matricielles suivantes sont possibles :
\[
  \left\{
  \begin{aligned}
    A_1X_1 + A_2X_2 &= B_1\\
    A_3X_1 + A_4X_2 &= B_2
  \end{aligned}
  \right.
\]
Si $A_1$ est régulière, alors on a
\[
  X_1 = A_1^{-1}\bigl[B_1 - A_2X_2\bigr].
\]
En substituant dans la deuxième équation, on obtient
\[
  \bigl[A_3A_1^{-1}A_2 - A_4\bigr]X_2 = \bigl[A_3A_1^{-1}B_1 - B_2\bigr].
\]
On peut ainsi calculer les éléments de $X_2$ et substituer dans l'expression
de $X_1$ pour avoir les éléments de $X_1$. On obtient donc ainsi toutes les
composantes de $X$.

Cette méthode peut être généralisée en décomposant la matrice en $9$, $16$,
$25$, \dots{} sous-matrices.

\subsubsection{Systèmes à coefficients complexes}

Soit le système matriciel
\[
  AZ = b
\]
avec
\[
  A = M + iN, \qquad b = c + id, \qquad Z = X + iY .
\]
Le système algébrique peut alors se décomposer en deux sous-systèmes de la
forme
\[
  \left\{
  \begin{aligned}
    MX - NY &= c\\
    NX + MY &= d
  \end{aligned}
  \right.
\]
On peut alors réunir ces deux sous-systèmes en un système unique, ou utiliser
l'approche du paragraphe précédent.

\section{Méthodes itératives pour les systèmes d'équations linéaires}

\subsection{Formules itératives}

Les méthodes directes telles que la méthode d'élimination de Gauss sont
appropriées pour les systèmes de taille moyenne. Pour les systèmes de grande
taille, il est conseillé d'utiliser les méthodes itératives qui sont plus
rapides, plus faciles à programmer et moins encombrantes. Cependant, ces
méthodes itératives sont basées sur le calcul des suites infinies qu'il faudra
tronquer à partir d'un certain rang. Pour le choix des valeurs initiales des
itérations, on fixe généralement les inconnues $x_i = 1$ ou $x_i = 0$ pour
$i = 1$ jusqu'à $n$.

Le principe des méthodes itératives consiste à exprimer l'inconnue $x_i$ de la
ligne $i$ en fonction des autres inconnues ; c'est-à-dire
\[
  x_i = \frac{b_i - \displaystyle\sum_{\substack{j=1\\ j\neq i}}^{n} a_{ij}x_j}{a_{ii}} .
\]
Ceci suggère que les valeurs connues $x_1, x_2, x_3, \dots, x_{i-1}, x_{i+1},
\dots, x_n$ soient substituées pour trouver $x_i$. Au départ, il faut donc
donner des valeurs initiales aux inconnues $x_i$. Ces valeurs initiales sont
généralement fixées à $x_i = 1$ ou $x_i = 0$ pour $i$ allant de $1$ à $n$.

Il existe deux variantes de méthode itérative : celle de \textbf{Jacobi} et
celle de \textbf{Gauss-Seidel}.

La formule itérative de \textbf{Jacobi} est :
\[
  x_i^{(k+1)} = \frac{b_i - \displaystyle\sum_{\substack{j=1\\ j\neq i}}^{n} a_{ij}x_j^{(k)}}{a_{ii}}
\]
où $x_i^{(k)}$ sont les valeurs des inconnues après $k$ itérations.

Dans la formule de \textbf{Gauss-Seidel}, les valeurs $x_j^{(k+1)}$ pour
$j \le i-1$ sont utilisées aussitôt qu'elles sont déterminées pour calculer
$x_i^{(k+1)}$. La formule itérative de \textbf{Gauss-Seidel} se met alors sous
la forme :
\[
  x_i^{(k+1)} = \frac{b_i - \displaystyle\sum_{j=1}^{i-1} a_{ij}x_j^{(k+1)}
                 - \displaystyle\sum_{j=i+1}^{n} a_{ij}x_j^{(k)}}{a_{ii}} .
\]
Notons que les deux formules sont importantes car il existe des systèmes
d'équations pour lesquels la méthode de Jacobi converge tandis que la méthode
de Gauss-Seidel ne converge pas. Dans d'autres cas c'est la méthode de
Gauss-Seidel qui converge et celle de Jacobi ne converge pas.

\subsection{Conditions d'arrêt}

Pour la condition d'arrêt de l'itération, il y a plusieurs méthodes. Elles
consistent à arrêter les itérations lorsque toutes les quantités
\[
  d_i = \abs{x_i^{(k+1)} - x_i^{(k)}} < \varepsilon
\]
ou aussi
\[
  d_i = \abs{\frac{x_i^{(k+1)} - x_i^{(k)}}{x_i^{(k+1)}}} < \varepsilon .
\]
Cependant, le meilleur test de convergence est d'utiliser la condition
\[
  d = \frac{\displaystyle\sum_{i=1}^{n}\abs{x_i^{(k+1)} - x_i^{(k)}}}
           {\displaystyle\sum_{i=1}^{n}\abs{x_i^{(k+1)}}} < \varepsilon .
\]

\begin{exercice}
Écrire les algorithmes des méthodes de Jacobi et de Gauss-Seidel avec
imposition de la condition d'arrêt.
\end{exercice}
